Deployed language models rarely run without a system prompt. Product teams add safety policies, personas, output-format rules and style constraints, and these prompts stay in context for every request, including requests they were never written for. We ask whether such prompts change accuracy, output length and refusal behaviour on unrelated tasks. We run a controlled experiment with 15 system prompts (three controls, plus four categories with three wording variants each) on five benchmarks (\gsm, \arc, a sensitive-subject slice of MMLU, \tqa and \drop) and three open instruction-tuned models, for 67,500 greedy generations and a full sampling replicate of the same size. We measure every effect as a paired, item-level difference from the minimal prompt ``You are a helpful assistant.'' The results contradict common expectations. Safety prompts leave accuracy unchanged ($+0.2$\pp, 95\% CI $[-0.4, 0.9]$), raise \tqa accuracy by 3.7\pp, and increase hedging without causing refusals. Persona prompts double output length ($1.99\times$) and lower accuracy by 3.1\pp. Format prompts lower answer correctness by 2.5\pp but raise strict-parser accuracy by 19.6\pp, so their value depends on what one measures. Behavioural constraints such as ``keep it under 50 words'' are the most harmful category ($-9.0$\pp on average, down to $-61$\pp in one cell), and they never help. All conclusions hold under temperature-0.7 sampling (cell-level $r=0.95$). For practitioners, length and style constraints in a system prompt deserve far more scrutiny than safety text.
